\subsection{COMPLETION OF A TOPOLOGICAL GROUP}

Let G be a Hausdorff topological group. The uniform space \(G_{d}\) may be considered as a dense subspace of its completion \(\hat{G}_{d}\) . We shall investigate whether G can be considered as a dense subgroup of a complete Hausdorff group \(G'\) . If so, then the uniform space \(G_{d}'\) must be isomorphic to \(\hat{G}_{d}\) (Chapter II, § 3, no. 6, Corollary to Theorem 2), and we must therefore be able to define on \(\hat{G}_{d}\) a topological group structure which induces the given topological group structure on G. Consequently we have to examine: 1) Whether we can extend by continuity the functions xy and \(x^{-1}\) to \(\hat{G}_{d} \times \hat{G}_{d}\) and \(\hat{G}_{d}\) respectively; 2) whether the functions thus extended do indeed define a group structure on \(\hat{G}_{d}\) (they will then necessarily define a topological group structure on \(\hat{G}_{d}\) inducing the given structure on G). We have next to establish that 3) when the preceding operations are possible, the topological group which they define is complete. Finally, we shall see that 4) if there is a complete group satisfying the given conditions, then it is unique up to isomorphism.

\begin{enumerate}
\def\labelenumi{\arabic{enumi})}
\tightlist
\item
  Extension of \(xy\) and \(x^{-1}\) by continuity. Since the functions \(xy\) and \(x^{-1}\) are not in general uniformly continuous, we cannot apply the theorem of extension of uniformly continuous functions (Chapter II, § 3, no. 6, Theorem 2). Nevertheless, we can extend \(xy\) , by virtue of Proposition 11 of Chapter II, § 3, no. 6 and the following Proposition:
\end{enumerate}

PROPOSITION 6. Let \(\mathfrak{F}\) and \(\mathfrak{G}\) be two Cauchy filters on \(\mathbf{G}_d\) . Then the image of the filter \(\mathfrak{F} \times \mathfrak{G}\) under the mapping \((x, y) \to xy\) is a Cauchy filter base on \(\mathbf{G}_d\) . Let us evaluate the ``proximity'' of \(xy\) and \(x'y'\) in \(\mathbf{G}_d\) ; in other words, let us form the product \((x'y')(xy)^{-1} = x'y'y^{-1}x^{-1}\) . For each \(a \in \mathbf{G}\) , we can also write \((x'y')(xy)^{-1} = (x'a^{-1})(ay'y^{-1}a^{-1})(ax^{-1})\) . We shall see that by suitable choice of \(a\) , each of the three factors of this product is very small whenever the pairs \((x, y)\) and \((x', y')\) belong to a sufficiently small set of \(\mathfrak{F} \times \mathfrak{G}\) . Let V be any neighbourhood of \(e\) in G; then there is a \(\mathbf{V}_d\) -small set A ∈ \(\mathfrak{F}\) . Choose \(a\) in A, then if \(x\) and \(x'\) are any two points of A, we have \(x'a^{-1} \in V\) and \(ax^{-1} \in V\) . On the other hand, the relation \(ay'y^{-1}a^{-1} \in V\) is equivalent to

\[
y ^ {\prime} y ^ {- 1} \in a ^ {- 1} \mathrm{V}a = \mathrm{W},
\]

and since W is a neighbourhood of e, there is a \(W_{d}\) -small set \(B \in \mathfrak{G}\) . Hence for all \((x, y)\) and \((x', y')\) in A × B, we have \((x'y')(xy)^{-1} \in V^{3}\) , and this completes the proof.

In order that \(x^{-1}\) can be extended by continuity to \(\hat{G}_d\) it is necessary and sufficient that the image, under the symmetry \(x \to x^{-1}\) , of a Cauchy filter on \(G_d\) is a Cauchy filter on \(G_d\) (Chapter II, § 3, no. 6, Proposition 11). There are examples of topological groups in which this condition is not satisfied (cf.~Chapter X, § 3, Exercise 16); we shall suppose that it is satisfied for the remainder of this proof.

\begin{enumerate}
\def\labelenumi{\arabic{enumi})}
\setcounter{enumi}{1}
\item
  The extended functions \(xy\) and \(x^{-1}\) define a group structure on \(\hat{G}_d\) . For if we apply the principle of extension of identities (Chapter I, § 8, no. 1, Proposition 2, Corollary 1) to the functions \(x(yz)\) and \((xy)z\) , defined on \(\hat{G}_d \times \hat{G}_d \times \hat{G}_d\) and equal on the dense subspace \(G_d \times G_d \times G_d\) , we see that the law of composition \((x, y) \to xy\) is associative on \(\hat{G}_d\) . For the same reason, the functions \(x, ex, xe\) are identical on \(\hat{G}_d\) , and the functions \(e, xx^{-1}, x^{-1}x\) are identical on \(\hat{G}_d\) .
\item
  The topological group \(\hat{\mathbf{G}}_d\) is complete. Let \(\mathfrak{U}_d\) be its right uniformity, and let \(\mathfrak{U}\) be the uniformity on \(\hat{\mathbf{G}}_d\) obtained by completing the right uniformity of \(\mathbf{G}\) . Then \(\mathfrak{U}\) and \(\mathfrak{U}_d\) induce the same uniformity on \(\mathbf{G}\) , and therefore every Cauchy filter base \(\mathfrak{B}\) on \(\mathbf{G}\) with respect to \(\mathfrak{U}_d\) is also a Cauchy filter base with respect to \(\mathfrak{U}\) . Now \(\mathfrak{B}\) converges in \(\hat{\mathbf{G}}_d\) , since \(\mathcal{U}\) is a complete uniformity; and since \(\mathcal{U}\) and \(\mathcal{U}_d\) induce the same topology on \(\hat{\mathbf{G}}_d\) , it follows (Chapter II, § 3, no. 4, Proposition 9) that \(\mathcal{U}_d\) is a complete uniformity. This conclusion also shows that \(\mathcal{U}\) and \(\mathcal{U}_d\) coincide (Chapter II, § 3, no. 6, Corollary to Theorem 2).
\item
  Uniqueness. This follows from Proposition 5 of no. 3.
\end{enumerate}

To sum up, we have proved the following theorem:

THEOREM 1. A Hausdorff topological group G is isomorphic to a dense subgroup of a complete group \(\hat{G}\) if and only if the image, under the symmetry \(x \to x^{-1}\) , of a Cauchy filter with respect to the right uniformity of G is a Cauchy filter with respect to this uniformity. The complete group \(\hat{G}\) (which is called the completion of G) is then unique (up to isomorphism).

PROPOSITION 7. Let G be a Hausdorff topological group which has a completion \(\hat{G}\) . Then the closures in \(\hat{G}\) of the neighbourhoods of the identity element in G form a fundamental system of neighbourhoods of the identity element in \(\hat{G}\) .

Since \(\hat{G}\) is regular, every neighbourhood of the identity element in \(\hat{G}\) contains the closure V of an open neighbourhood U of e in \(\hat{G}\) , and V is also the closure of the trace of U on G.

Let G be a group which is not necessarily Hausdorff; let \(N = \overline{e}\) , and let \(G' = G/N\) be the Hausdorff group associated with G (§ 2, no. 6). If \(G'\) has a completion \(\hat{G}'\) , this completion is called the Hausdorff completion of G and is denoted by \(\hat{G}\) ; \(\hat{G}_{d}'\) (resp. \(\hat{G}_{s}'\) ) is then the Hausdorff completion (Chapter II, § 3, no. 7) of the uniform space \(G_{d}\) (resp. \(G_{s}\) ).

PROPOSITION 8. Let G be a topological group which has a Hausdorff completion \(\hat{G}'\) . Then every continuous homomorphism \(u\) of G into a complete Hausdorff group H can be uniquely factorized into \(u = v \circ \varphi\) , where \(v\) is a continuous homomorphism of \(\hat{G}'\) into H and \(\varphi\) is the canonical mapping of G into \(\hat{G}'\) (the composition of the canonical injection of G' into \(\hat{G}'\) and the canonical homomorphism \(\psi\) of G onto G/N = G').

Since the kernel of \(u\) is closed and contains \(e\) , it contains \(N\) , and hence \(u\) can be written as \(u = w \circ \psi\) , where \(w\) is a continuous homomorphism of \(G'\) into \(H\) ; now apply Proposition 5 of no. 3 to \(w\) .

